        	\subsection{The Morse index of half-bubbles}\label{subs:morsehalf}
        	
        	The discussion presented in the previous section directly applies, as a special case, to the study of the Jacobi operator
        	that is associated to the symmetrized complete minimal hypersurface that is obtained by reflecting a half-bubble. 
        	
        	Given a half-bubble $\Sigma$ we define, in analogy with Definition \ref{def:index} above, its Morse index considering the Jacobi form
\[        	
Q_{|A|^2}(u,u):=\int_{\Sigma}(|\nabla u|^2-|A|^2u^2)\,d\h^n.
\] 
More precisely, we give the following:
        	
        	\begin{definition} In the setting above, we define
        		$\operatorname{index} (\Sigma)$ to be the largest dimension of a linear subspace of $C^{\infty}_c(\Sigma)$ where $Q_{|A|^2}$ is negative definite. 
        	\end{definition}	
        	\noindent Notice that here the support of $\varphi\in C^{\infty}_c(\Sigma)$ can intersect $\partial\Sigma$ so that we are imposing no condition along $\partial\Sigma$, and $\operatorname{index} (\Sigma)$ is the (standard) Morse index of $\Sigma$ as a free boundary minimal hypersurface.
        	
        	\begin{lem}\label{lem:II1}
        A two-dimensional half-bubble has finite Morse index and this equals $\operatorname{Ind}_{{\jepPubScript{E}}}(Q_{|\check{A}|^2})$ where $\check{\Sigma}$ is the double of $\Sigma$.
         The same conclusion holds true in all dimensions for any half-bubble $\Sigma$ having Euclidean volume growth, meaning that there exists a constant $C=C(\Sigma)$ such that $\h^n(\Sigma\cap B_R(0))\leqslant C R^n$ for all $R\geqslant 0$.
        	\end{lem}
        
           \begin{proof}
        	The argument for the first part ($n=2$) goes as follows: by definition $\Sigma$ has finite total curvature, so its double $\check{\Sigma}$ will also have finite total curvature, hence finite Morse index (again by \cite{FC85}); thus Lemma \ref{lem:ineq} implies that the corresponding \textsl{even} Morse index (that is to say: the Morse index on even test functions, as per Definition \ref{def:index}) is also finite. Now, going back to the definitions it is clear that by restriction $\operatorname{index} (\Sigma)\geqslant \operatorname{Ind}_{{\jepPubScript{E}}}(Q_{|\check{A}|^2})$.
        	
        	For the converse inequality, we argue by contradiction as follows. If it were $\operatorname{index} (\Sigma)> \operatorname{Ind}_{{\jepPubScript{E}}}(Q_{|\check{A}|^2})$, we could consider a basis $\left\{\varphi^1,\ldots, \varphi^I\right\}$ for a subspace of $C^{\infty}_c(\Sigma)$, having dimension equal to $I:=\operatorname{Ind}_{{\jepPubScript{E}}}(Q_{|\check{A}|^2})+1$, where $Q_{|A|^2}$ is negative definite. 
        			If we extend each function by even symmetry, we obtain a family of even functions $\left\{\check{\varphi}^1,\ldots \check{\varphi}^I\right\}$ where the Jacobi form of $\check{\Sigma}$ is negative definite; each function is smooth away from $\partial\Sigma$  and we can consider for each $i\in\left\{1,\ldots, I\right\}$ an approximating sequence of smooth functions $\left\{\check{\varphi}^i_k\right\}$ that also have compact support, and such that $\check{\varphi}^i_k\to\check{\varphi}^i$ uniformly as one lets $k\to\infty$.
        	 		 We claim that for any $k$ large enough the family $\left\{\check{\varphi}^1_k,\ldots \check{\varphi}^I_k\right\}$ is linearly independent, which would then violate the definition of $\operatorname{Ind}_{{\jepPubScript{E}}}(Q_{|\check{A}|^2})$ and conclude the proof.
        			 If the claim were false, we could write down (for every $k\in\mathbb{N}$) a linear equation of the form
        		 \[
        		 \sum_i a^i_k \check{\varphi}^i_k=0, \ \text{where} \ \sum_i (a^i_k)^2=1,
        		 \]
        		whence passing to the limit for $k\to\infty$, we end up finding a (non-trivial) linear relation involving 
        	$\left\{\varphi^1,\ldots, \varphi^I\right\}$, that is impossible since this family of functions was chosen to be a basis.
        	 In higher dimensions, one can follow the same argument modulo invoking the main theorem of J. Tysk in \cite{T89}.
        	\end{proof}
       
        
        	
        	We can instead consider the Morse index of $\Sigma$ with Dirichlet boundary conditions:
        	
        	\begin{definition} In the setting above, we define
        		$\operatorname{index}_{\jepPubBullet} (\Sigma)$ to be the largest dimension of a linear subspace of $C^{\infty}_c(\mathring{\Sigma})$ where $Q_{|A|^2}$ is negative definite.
        	\end{definition}
        	
        	Here $\mathring{\Sigma}:=\Sigma\smallsetminus\partial\Sigma$ is the interior of $\Sigma$, and so this is the standard notion of Morse index for minimal surfaces with respect to variations that fix the boundary. Following the same conceptual scheme as above, we have the following ancillary result:
        	
        	\begin{lem}\label{lem:II2}
        		A two-dimensional half-bubble has finite Morse index \emph{with Dirichlet boundary conditions} and this equals $\operatorname{Ind}_{{\jepPubScript{O}}}(Q_{|\check{A}|^2})$ where $\check{\Sigma}$ is the double of $\Sigma$.
        		The same conclusion holds true in all dimensions for any half-bubble $\Sigma$ having Euclidean volume growth, meaning that there exists a constant $C=C(\Sigma)$ such that $\h^n(\Sigma\cap B_R(0))\leqslant C R^n$ for all $R\geqslant 0$.
        	\end{lem}	
        	
        	\begin{proof}
        	Arguing as for the lemma above, let us focus on $n=2$. When $\Sigma$ is a two-dimensional half-bubble, the finiteness of $\operatorname{index}_{\jepPubBullet} (\Sigma)\leqslant \operatorname{index}(\Sigma)$ follows straight from Lemma \ref{lem:II1}. That being said, it is clear that we also have the inequality
        	$\operatorname{index}_{\jepPubBullet} (\Sigma)\leqslant \operatorname{Ind}_{{\jepPubScript{O}}}(Q_{|\check{A}|^2})$ by odd extension. The converse inequality
        	relies instead on the fact that odd functions must vanish along $\partial\Sigma=\operatorname{Fix}(\tau)$ so that one can consider a linearly independent set of $\operatorname{Ind}_{{\jepPubScript{O}}}(Q_{|\check{A}|^2})$ elements spanning a subspace on which $Q_{|\check{A}|^2}<0$, restrict each such function to $\Sigma$ and then truncate using a compactly supported cutoff function. We omit the standard details.
        		\end{proof}
        	
        Based on Lemma \ref{lem:II1} and \ref{lem:II2} one can rephrase the inequality $\operatorname{index}_{\jepPubBullet} (\Sigma)\leqslant \operatorname{index}(\Sigma)$ as
        	\begin{equation}\label{eq:EindexvsOindex}
        	\operatorname{Ind}_{{\jepPubScript{E}}}(Q_{|\check{A}|^2})\geqslant  \operatorname{Ind}_{{\jepPubScript{O}}}(Q_{|\check{A}|^2}),
        	\end{equation}
        	hence, using the second part of Lemma \ref{lem:specadd}, we get the estimate
        	\begin{equation}\label{eq:fundineq}
        	2 \operatorname{Ind}_{{\jepPubScript{O}}}(Q_{|\check{A}|^2})\leqslant \operatorname{Ind}(Q_{|\check{A}|^2})\leqslant 2 \operatorname{Ind}_{{\jepPubScript{E}}}(Q_{|\check{A}|^2}).
        	\end{equation}
        	From there, we can derive some interesting characterizations for the cases when the Morse index of $\check{\Sigma}$ (which is, in our notation, $	\operatorname{Ind}(Q_{|\check{A}|^2})$) equals 0, 1 or 2:
        	
        	\begin{cor}\label{cor:lowindex}
        		In the setting above, we have 
        		\begin{align*}
        		&\text{(1)}& 		\operatorname{Ind}(Q_{|\check{A}|^2})=0 \ \iff \operatorname{Ind}_{{\jepPubScript{E}}}(Q_{|\check{A}|^2})=0 \ \text{and} \ \operatorname{Ind}_{{\jepPubScript{O}}}(Q_{|\check{A}|^2})=0; \\
        		&\text{(2)}& 		\operatorname{Ind}(Q_{|\check{A}|^2})=1 \ \iff \operatorname{Ind}_{{\jepPubScript{E}}}(Q_{|\check{A}|^2})=1 \ \text{and} \ \operatorname{Ind}_{{\jepPubScript{O}}}(Q_{|\check{A}|^2})=0;\\
        		&\text{(3)}&		\operatorname{Ind}(Q_{|\check{A}|^2})=2 \ \iff 
        				\begin{cases}\operatorname{Ind}_{{\jepPubScript{E}}}(Q_{|\check{A}|^2})=2 \ \text{and} \ \operatorname{Ind}_{{\jepPubScript{O}}}(Q_{|\check{A}|^2})=0 \ \ \textbf{or} \\ 
        				 \operatorname{Ind}_{{\jepPubScript{E}}}(Q_{|\check{A}|^2})=1 \ \text{and} \ \operatorname{Ind}_{{\jepPubScript{O}}}(Q_{|\check{A}|^2})=1.
        				\end{cases}
        				\end{align*}
        	\end{cor}
        	
        	So far our discussion is applicable to hypersurfaces of dimension $n\geqslant 2$, and we shall now specify them to the special case $n=2$, where some interesting consequences can be drawn.
        	Indeed, we can turn classification results for complete minimal surfaces in $\R^{3}$ into classification results for free boundary minimal hypersurfaces contained in a half-space. The first one is a Bernstein-type theorem for half-bubbles:
        	
        	\begin{cor}\label{cor:stable}
        		A two-dimensional, stable half-bubble is isometric to a half-plane.	
        	\end{cor}	
        	
        	\begin{proof}
        		It suffices to combine (1) of Corollary \ref{cor:lowindex}, inequality \eqref{eq:EindexvsOindex} and the characterization of stable minimal surfaces in $\R^{3}$ (see \cite{dCP79,FCS80,Pog81}).
        	\end{proof}
        	
        	\begin{rmk}\label{rmk:higher-stable}
        	In fact, the stability estimates by Schoen-Simon allow to obtain a higher-dimensional counterpart of the previous result: \textsl{Let $\Sigma^n\subset\R^{n+1}$, $2\leqslant n\leqslant 6$ be a half-bubble. If $\Sigma$ is stable and has Euclidean volume growth, then $\Sigma$ is a hyperplane.}
        	\end{rmk}
        		
        	
        	
        	
        	Similarly, we get a simple result in the index one case:	
        	
        	\begin{cor}\label{cor:indexone}
        		A two-dimensional, index one half-bubble is isometric to a half-catenoid.	
        	\end{cor}	
        	
        	\begin{proof}
        		It suffices to combine inequality \eqref{eq:EindexvsOindex}, parts (2) and (3) of Corollary \ref{cor:lowindex}, the characterization of index one minimal surfaces in $\R^3$, see \cite{LR89}, and the non-existence result for index two complete minimal surfaces in $\R^3$, see \cite{CM14}.
        	\end{proof}
        	
        	In fact, inspired by recent work of Chodosh-Maximo \cite{CM14} we may wonder whether there exist two-dimensional half-bubbles whose Morse index equals two.
        	Obviously, a negative result would follow by proving that a half-bubble of Morse index 2 has vanishing Morse index with Dirichlet boundary conditions. Yet, this does not directly follow from the results above. 	
        	
